\par\medskip\noindent
\textbf{\textit{Reminder 4.10.}}---
\textit{Let \(k\) be an algebraically closed field, let \(G\) be a
reductive \(k\)-group, let \(T\) be a maximal torus of \(G\), and let
\(\alpha\) be a root of \(G\) with respect to \(T\).  There exists a
monomorphism}
\[
  p\colon\mathbb G_{a,k}\longrightarrow G
\]
\textit{normalized by \(T\) with multiplier \(\alpha\).}

See the \emph{Bible}, §13.1, Theorem 1.
\SGARefTarget{sga3:xix:reminder:4:10}

\par\medskip\noindent
\textbf{\SGAInlineRefTarget{sga3:xix:paragraph:4-11}4.11.}---
We finish this section with a technical result which will be useful
to us.  Let \(S\) be a scheme, and let \(L\) be an invertible
\(\mathcal O_S\)-module.  Let \(q>0\) be an integer such that
\(x\longmapsto x^q\) is an endomorphism of the \(S\)-group
\(\mathbb G_{a,S}\).  If \(S\ne\varnothing\), then \(q=1\), or
\(q=p^n\), where \(p\) is a prime number which vanishes on \(S\).
This follows at once from the elementary fact that the greatest
common divisor of the binomial coefficients
\(\binom qi\), for \(i\ne0,q\), is \(p\) if \(q=p^n\) with \(p\)
prime, and is \(1\) otherwise.

The morphism defined by taking the \(q\)-th power,
\[
  L\longrightarrow L^{\otimes q},
\]
is a morphism of sheaves of abelian groups.  By base change it
defines a morphism of \(S\)-group schemes
\[
  \mathbb W(L)\longrightarrow\mathbb W(L^{\otimes q}).
\]
In particular, if \(L'\) is another invertible module and an
\(\mathcal O_S\)-linear morphism
\[
  h\colon L^{\otimes q}\longrightarrow L'
\]
is given, one obtains a morphism of \(S\)-group schemes
\[
  \mathbb W(L)\longrightarrow\mathbb W(L'),
  \qquad
  x\longmapsto h(x^q).
\]
\SGARefTarget{sga3:xix:paragraph:4:11}

With this notation, we have the following result.

\par\medskip\noindent
\textbf{\textit{\SGARefLink{sga3:xix:proposition:4:12}{Proposition 4.12}.}}---
\textit{Let \(S\) be a scheme, let \(T\) \textup{(resp. \(T'\))} be
an \(S\)-torus, let \(L\) \textup{(resp. \(L'\))} be an invertible
\(\mathcal O_S\)-module, and let \(\alpha\)
\textup{(resp. \(\alpha'\))} be a character of \(T\)
\textup{(resp. \(T'\))}.  Let \(f\colon T\to T'\) be a group
morphism and let
\[
  g\colon\mathbb W(L)\longrightarrow\mathbb W(L')
\]
be a morphism of \(S\)-schemes \textup{(not necessarily a morphism
of \(S\)-groups)} satisfying}
\[
  \tag{\(\star\)}\label{sga3:xix:paragraph:4-11:equation:star}
  g\bigl(\alpha(t)x\bigr)
  =
  \alpha'\bigl(f(t)\bigr)g(x)
\]
\textit{for every \(S'\to S\), every \(x\in\mathbb W(L)(S')\), and
every \(t\in T(S')\).  Let \(s_0\in S\) be such that
\(\alpha_{\bar s_0}\ne0\).}
\footnote{Editor's note: the phrases ``consider the semidirect
product \(T\mathbin{\cdot_\alpha}\mathbb W(L)\), respectively
\(T'\mathbin{\cdot_{\alpha'}}\mathbb W(L')\)'' have been deleted.
On the other hand, in part \textup{\SGARefLink{sga3:xix:paragraph:4:11:item:b}{(b)}} the hypothesis that \(g\) be
a group morphism, together with \textup{\((\star)\)}, is equivalent
to requiring
\[
  (t,x)\longmapsto\bigl(f(t),g(x)\bigr)
\]
to be a group morphism from
\(T\mathbin{\cdot_\alpha}\mathbb W(L)\) to
\(T'\mathbin{\cdot_{\alpha'}}\mathbb W(L')\).}

\begin{enumerate}
\item[\textup{(a)}]\SGAInlineRefTarget{sga3:xix:proposition:4:12:item:a}
\textit{Suppose that \(g\) sends the zero section to the zero section
and that, for every integer \(n>0\),}
\[
  (\alpha'\circ f)_{\bar s_0}\ne n\alpha_{\bar s_0}.
\]
\textit{Then \(g=0\) in a neighborhood of \(s_0\).}

\item[\textup{\SGAInlineRefTarget{sga3:xix:proposition:4-12:item:b}\SGAInlineRefTarget{sga3:xix:paragraph:4:11:item:b}(b)}]
\textit{Suppose that \(g\) is a group morphism and that
\(g_{\bar s_0}\ne0\).  There then exist an open subset \(U\subset S\)
containing \(s_0\) and an integer \(q>0\) such that
\(x\longmapsto x^q\) is an endomorphism of
\(\mathbb G_{a,U}\) and}
\[
  (\alpha'\circ f)_U=q\alpha_U.
\]
\end{enumerate}
\SGARefTarget{sga3:xix:proposition:4:12}

